\subsection{Cross-Mode Learning and Curve Shape}\label{sec:cross-mode-learning}

The paper predicts that at high $m$, the $a_k$ curve should be sigmoidal: an initial plateau (early responses come from different modes and are uninformative about each other), followed by decline once mode repetitions accumulate.
GPT-2 shows this pattern at $m=10$.
Qwen2.5-3B does not: it shows immediate exponential decay at all $m$.
This section investigates the discrepancy through systematic pairwise analysis.

\paragraph{Pairwise cross-mode surprise matrix.} For 15 modes, we compute a $15 \times 15$ matrix $M_{ij}$: the surprise reduction (in bits) when a response from mode $j$ is in the context and mode $i$ is the target.
Diagonal entries measure same-mode learning; off-diagonal entries measure cross-mode information transfer.
Each entry averages over 5 context--target samples (using different samples for context and target to avoid self-prediction inflation).

\begin{figure*}[p]
\centering
\begin{subfigure}{\textwidth}
    \centering
    \includegraphics[width=\textwidth]{pairwise_matrix_heatmap.png}
    \caption{Qwen2.5-3B: positive off-diagonal.}
    \label{fig:pairwise-qwen}
\end{subfigure}

\vspace{1.5em}

\begin{subfigure}{\textwidth}
    \centering
    \includegraphics[width=\textwidth]{gpt2/pairwise_matrix_heatmap.png}
    \caption{GPT-2: negative off-diagonal.}
    \label{fig:pairwise-gpt2}
\end{subfigure}
\caption{Pairwise cross-mode surprise reduction matrices.
Each cell $(i,j)$ shows how many bits of surprise reduction mode $i$ (target) receives from seeing mode $j$ (context).
Qwen shows diagonal dominance with pervasive positive off-diagonal (\crossmodeQwenOffMean{} bits mean); GPT-2 shows diagonal dominance with negative off-diagonal (\crossmodeGPTOffMean{} bits mean).}
\label{fig:pairwise-matrices}
\end{figure*}

\paragraph{Results.} (Cell standard deviations are the mean within-cell sampling std across the \crossmodeNumModes{} diagonal or \crossmodeNumOffPairs{} off-diagonal matrix cells, each estimated from 5 context-target sample pairs.)
\begin{itemize}[nosep]
    \item \textbf{Qwen2.5-3B}: diagonal mean = $\crossmodeQwenDiagMean \pm \crossmodeQwenDiagStd$ bits, off-diagonal mean = $\crossmodeQwenOffMean \pm \crossmodeQwenOffStd$ bits, \crossmodeQwenOffPctPos\% of off-diagonal entries positive.
    \item \textbf{GPT-2}: diagonal mean = $\crossmodeGPTDiagMean \pm \crossmodeGPTDiagStd$ bits, off-diagonal mean = $\crossmodeGPTOffMean$ bits, only \crossmodeGPTOffPctPos\% of off-diagonal entries positive.
\end{itemize}

The sign of the off-diagonal determines the curve shape (Figure~\ref{fig:pairwise-matrices}):
\begin{itemize}[nosep]
    \item \textbf{Positive off-diagonal (Qwen)}: Every response helps predict every other response through format calibration and topic priming.
The $a_k$ curve drops from position 1, producing exponential decay.
    \item \textbf{Negative off-diagonal (GPT-2)}: Cross-mode responses actively interfere with predictions.
At $m=10$, the expected first-step net drop under the additive independent-mode model is only $\crossmodeGPTFirstStepDrop$ bits (the \crossmodeGPTDiagGainMten{}-bit diagonal gain from the $1/m$ same-mode chance is largely offset by the $\crossmodeGPTCrossDamageMten$-bit cross-mode damage from the $(m{-}1)/m$ different-mode chance), producing a flat plateau until same-mode repetitions accumulate, the predicted sigmoid.
\end{itemize}

\paragraph{From pairwise matrix to full curve.} The sign-of-off-diagonal calculation above recovers the qualitative shape difference between Qwen and GPT-2 from pairwise measurements alone.
Extending to the full $a_k$ curve at $m>1$ requires modeling information overlap across repeated encounters --- structure that pairwise (single-encounter) measurements do not capture.
Additive and fractional decompositions of the pairwise matrix both fail in instructive ways but neither is load-bearing for our thesis; we omit the details.

\paragraph{Pairwise asymmetry.} The mean absolute asymmetry $|M_{ij} - M_{ji}|$ is \crossmodeAsymmetryMean{} bits ($\crossmodeAsymmetryRatio\times$ the off-diagonal mean), falsifying the hypothesis that pairwise surprise reduction is symmetric (Figure~\ref{fig:pairwise-symmetry}).
This reflects different mechanisms: technical/structured modes (code, JSON, math) are highly informative as context (format calibration) but do not benefit proportionally from other modes' context.

\begin{figure*}[t]
\centering
\includegraphics[width=0.5\textwidth]{pairwise_symmetry.png}
\caption{Pairwise symmetry scatter for Qwen2.5-3B: each point is one $(i,j)$ pair, plotting $M_{ij}$ vs.\ $M_{ji}$.
Under symmetric mutual information, points would lie on $y = x$.
The best-fit slope is \scalingSymSlopeQwen{} ($R^2 = \scalingSymRsqQwen$), and many pairs have opposite signs: mode $i$ can reduce surprise for mode $j$ while $j$ \emph{increases} surprise for $i$.}
\label{fig:pairwise-symmetry}
\end{figure*}

\paragraph{Connections to information-theoretic properties.} Two information-theoretic properties are relevant to interpreting the pairwise matrix:

\begin{enumerate}[nosep]
    \item \textbf{Non-negative conditioning} ($H(X) \geq H(X \mid Y)$): Under the true distribution, conditioning can never increase entropy on average.
For cross-entropies this is not guaranteed (the KL term can increase), but a well-calibrated $\theta$ should show mostly non-negative off-diagonal entries.
    \item \textbf{Symmetry of mutual information} ($I(X;Y) = I(Y;X)$): Under the true joint, row means (how informative mode $i$ is as context) should correlate with column means (how much mode $i$ benefits from context).
\end{enumerate}

\begin{figure*}[t]
\centering
\begin{subfigure}[t]{0.48\textwidth}
    \includegraphics[width=\textwidth]{pairwise_row_vs_col.png}
    \caption{Qwen2.5-3B: slope = \scalingRCSlopeQwen{}, $R^2 = \scalingRCRsqQwen$.}
    \label{fig:rowcol-qwen}
\end{subfigure}
\hfill
\begin{subfigure}[t]{0.48\textwidth}
    \includegraphics[width=\textwidth]{gpt2/pairwise_row_vs_col.png}
    \caption{GPT-2: slope = \scalingRCSlopeGPT{}, $R^2 = \scalingRCRsqGPT$.}
    \label{fig:rowcol-gpt2}
\end{subfigure}
\caption{Row mean (how informative as context) vs.\ column mean (how much benefit from context) for each mode.
Qwen shows tighter correlation closer to the identity, consistent with approximate symmetry of mutual information.
GPT-2 shows widespread violations of non-negative conditioning (modes in the negative-row-mean region).}
\label{fig:row-vs-col}
\end{figure*}

Figure~\ref{fig:row-vs-col} compares the two models.
Qwen shows a tighter row--column relationship (slope = \scalingRCSlopeQwen{}, $R^2 = \scalingRCRsqQwen$) with only \crossmodeQwenNegColCount{} of \crossmodeNumModes{} modes having a negative column mean.
GPT-2 shows widespread violations of non-negative conditioning (\crossmodeGPTNegColCount{} modes with negative column means) and a barely-existent row--column relationship (slope = \scalingRCSlopeGPT{}, $R^2 = \scalingRCRsqGPT$).
With only two base models tested, we have suggestive but not conclusive evidence that stronger models better satisfy these information-theoretic properties.

\paragraph{Token-level mechanisms.} Per-token attribution analysis reveals two mechanisms for cross-mode surprise reduction:
\begin{itemize}[nosep]
    \item \textbf{Format calibration} (front-loaded): Seeing any structured response recalibrates the model's expectations about what follows ``Response B:''.
This effect concentrates in the first few tokens of the target response; the per-pair magnitude varies widely (positive for most cross-mode pairs, negative when the context format actively conflicts with the target).
    \item \textbf{Topic/vocabulary priming} (distributed): All responses are about rain, so rain-related vocabulary gets primed cross-mode.
This contribution is spread across the target response rather than concentrated at the start.
\end{itemize}

\paragraph{Discussion.} The paper's sigmoid prediction assumes modes are approximately independent.
For capable models, this fails: cross-mode learning (format calibration, topic priming) creates positive information transfer between distinct modes.
The comparison between GPT-2 and Qwen is suggestive: the stronger model shows behavior closer to information-theoretic ideals.
If this pattern holds across more models, it would mean the metric's theoretical properties are better satisfied as $\theta$'s ICL capability improves, but confirming this requires testing additional base models of varying sizes.

